% Part: normal-modal-logic
% Chapter: axioms-systems
% Section: normal-logics

\documentclass[../../../include/open-logic-section]{subfiles}

\begin{document}

\olfileid{nml}{prf}{nor}

\olsection{المنطقات الجهية النظامية}

لا يسهل أن نجعل كل مجموعة من !!{formula}s الجهية مجموعة
!!{formula}s مشتقة من بديهيات؛ بل نطلب شروطًا تضبط المنطق.
فما يصح أن !!{derive} في منطق القضايا الكلاسيكي ينبغي أن يبقى
!!{derivable}، وإن اتسعت !!{formula}s فدخل فيها
\iftag{prvBox}{$\Box$\iftag{prvDiamond}{
    و~}{}}{}\iftag{prvDiamond}{$\Diamond$}{}.
ولهذا نشترط جميع أمثلة القضايا الصادقة صدقًا تحليليًا، ونشترط
الانغلاق تحت القياس الاستثنائي.

\begin{defn}
  \emph{المنطق الجهوي} هو مجموعة~$\Sigma$ من !!{formula}s الجهية بحيث
  \begin{enumerate}
  \item تشتمل على جميع القضايا الصادقة صدقا تحليليا، و
  \item تكون مغلقة تحت الإحلال، أي إذا كان $!A \in \Sigma$ وكانت
    $!D_{\OLMathNumeral{1}}$، \dots، $!D_{\OLId{latin-ordinary}{n}}$ هي !!{formula}s، فعندئذ
    \[
    \SSubst{!A}{\subst{!D_{\OLMathNumeral{1}}}{{\OLId{latin-ordinary}{p}}_{\OLMathNumeral{1}}}, \dots, \subst{!D_{\OLId{latin-ordinary}{n}}}{{\OLId{latin-ordinary}{p}}_{\OLId{latin-ordinary}{n}}}} \in \Sigma,
    \]
  \item تكون مغلقة تحت \emph{القياس الاستثنائي}، أي إذا كان $!A$
    و$!A \lif !B \in \Sigma$، فعندئذ $!B \in \Sigma$.
  \end{enumerate}
\end{defn}

ثم إن ملاءمة الدلالة العلاقية تقتضي اشتمال المنطق على جميع
!!{formula}s الصالحة في النماذج الجهية. ويكفي لذلك أن تكون
كل أمثلة~\Ax{K} \iftag{prvDiamond}{و~\Dual{}}{} !!{derivable}،
وأن يكون، كلما كانت !!a{formula}~$!A$ !!{derivable}،
كذلك~$\Box !A$. فهذا هو المنطق الجهوي \emph{النظامي}.
ولا ننفي وجود منطقات غير نظامية، وإنما نقول إن النماذج العلاقية
المعتادة لا تكفي لدلالتها.

\begin{defn}
  يكون المنطق الجهوي $\Sigma$ \emph{نظاميا} إذا اشتمل على
  \begin{align*}
    \tag{\Ax{K}} & \Box({\OLId{latin-ordinary}{p}} \lif {\OLId{latin-ordinary}{q}}) \lif (\Box {\OLId{latin-ordinary}{p}} \lif \Box {\OLId{latin-ordinary}{q}}),
    \iftag{prvDiamond}{\\
      \tag{\Dual} & \Diamond {\OLId{latin-ordinary}{p}} \liff \lnot\Box\lnot {\OLId{latin-ordinary}{p}}}{}
  \end{align*}
  وكان مغلقا تحت \emph{الإلزام}، أي إذا كان $!A \in
  \Sigma$، فعندئذ $\Box !A \in \Sigma$.
\end{defn}

وهنا فرق في مرتبة حفظ الصدق: اللزوم التحليلي يحفظ
\emph{الصدق في عالم} بعينه؛ وأما~\Nec{} فلا تحفظ إلا
\emph{الصدق في نموذج} كله. ومن ذلك يلزم حفظها للصلاحية
في إطار أو في فئة أطر.

\begin{prop}\ollabel{prop:rk}
  كل منطق جهوي نظامي مغلق تحت القاعدة \RK،
  \begin{prooftree}
    \AxiomC{$!A_{\OLMathNumeral{1}} \lif (!A_{\OLMathNumeral{2}} \lif \cdots (!A_{{\OLId{latin-ordinary}{n}}-{\OLMathNumeral{1}}} \lif !A_{\OLId{latin-ordinary}{n}})\cdots)$}
    \RightLabel{\RK}
    \UnaryInfC{$\Box!A_{\OLMathNumeral{1}} \lif (\Box!A_{\OLMathNumeral{2}} \lif \cdots (\Box!A_{{\OLId{latin-ordinary}{n}}-{\OLMathNumeral{1}}}
      \lif \Box!A_{\OLId{latin-ordinary}{n}})\cdots).$}
  \end{prooftree}
\end{prop}

\begin{proof}
  بالاستقراء على~${\OLId{latin-ordinary}{n}}$: إذا كان ${\OLId{latin-ordinary}{n}} = {\OLMathNumeral{1}}$، فالقاعدة ليست سوى \Nec، وكل
  منطق جهوي نظامي مغلق تحت \Nec.

  ولنفترض الآن أن النتيجة صحيحة لـ${\OLId{latin-ordinary}{n}}-{\OLMathNumeral{1}}$؛ ونبين صحتها لـ${\OLId{latin-ordinary}{n}}$.

  افترض أن
  \begin{align*}
  & !A_{\OLMathNumeral{1}} \lif (!A_{\OLMathNumeral{2}} \lif \cdots (!A_{{\OLId{latin-ordinary}{n}}-{\OLMathNumeral{1}}} \lif !A_{\OLId{latin-ordinary}{n}})\cdots) \in \Sigma
  \intertext{وبفرض الاستقراء، لدينا}
  & \Box!A_{\OLMathNumeral{1}} \lif (\Box!A_{\OLMathNumeral{2}} \lif \cdots \Box(!A_{{\OLId{latin-ordinary}{n}}-{\OLMathNumeral{1}}} \lif !A_{\OLId{latin-ordinary}{n}})\cdots)
  \in \Sigma
  \intertext{ولأن $\Sigma$ منطق جهوي نظامي، فإنه يشتمل على جميع أمثلة
    $\Ax{K}$، ومنها بوجه خاص}
  & \Box(!A_{{\OLId{latin-ordinary}{n}}-{\OLMathNumeral{1}}} \lif !A_{\OLId{latin-ordinary}{n}}) \lif (\Box!A_{{\OLId{latin-ordinary}{n}}-{\OLMathNumeral{1}}} \lif \Box!A_{\OLId{latin-ordinary}{n}}) \in \Sigma
  \intertext{وباستخدام القياس الاستثنائي وأمثلة مناسبة لقضايا صادقة صدقا تحليليا نحصل على}
  & \Box!A_{\OLMathNumeral{1}} \lif (\Box!A_{\OLMathNumeral{2}} \lif \cdots (\Box!A_{{\OLId{latin-ordinary}{n}}-{\OLMathNumeral{1}}}
  \lif \Box!A_{\OLId{latin-ordinary}{n}})\cdots) \in \Sigma.
  \end{align*}
\end{proof}

\begin{prop}\ollabel{prop:notDiamondBot}
  يشتمل كل منطق جهوي نظامي $\Sigma$ على~$\lnot\Diamond\lfalse$.
\end{prop}

\begin{prob}
  أثبت \olref[nml][prf][nor]{prop:notDiamondBot}.
\end{prob}

\begin{prop}
  لتكن $!A_{\OLMathNumeral{1}}$، \dots، $!A_{\OLId{latin-ordinary}{n}}$ هي !!{formula}s. عندئذ يوجد أصغر
  منطق جهوي نظامي $\Sigma$ يشتمل على جميع أمثلة $!A_{\OLMathNumeral{1}}$، \dots، $!A_{\OLId{latin-ordinary}{n}}$.
\end{prop}

\begin{proof}
  % تصحيح مصدر (0429-I1): وجود مجموعة الصيغ يثبت عدم خلو أسرة التقاطع لا التقاطع نفسه.
  خذ~$!A_{\OLMathNumeral{1}}$، \dots، $!A_{\OLId{latin-ordinary}{n}}$، واجعل~$\Sigma$ تقاطع المنطقات
  الجهية النظامية التي تشتمل على كل أمثلة~$!A_{\OLMathNumeral{1}}$، \dots، $!A_{\OLId{latin-ordinary}{n}}$.
  وهذه الأسرة التي نأخذ تقاطعها غير خالية؛ إذ إن~$\Frm[{\OLId{latin-looped}{L}}]$،
  مجموعة جميع !!{formula}s، واحد من تلك المنطقات.
  وكل شرط من شروط المنطق النظامي يبقى عند التقاطع؛ فيكون التقاطع
  منطقًا نظاميًا، ويكون أصغرها بالاحتواء.
\end{proof}

\begin{defn}
يسمى أصغر منطق جهوي نظامي يشتمل على $!A_{\OLMathNumeral{1}}$، \dots، $!A_{\OLId{latin-ordinary}{n}}$
\emph{نسقا جهيا}، ويرمز إليه بـ$\Log{K} !A_{\OLMathNumeral{1}} \dots
!A_{\OLId{latin-ordinary}{n}}$. ويرمز إلى أصغر منطق جهوي نظامي بـ~\Log{K}.
\end{defn}

\end{document}
